\subsection{A second independent quadratic annihilator}

The full canonical central cover in B.2.6 can be computed
without choosing an abstract Schur cover.  The computation
below uses the marked source itself at every step and,
together with B.2.7, produces an actual full profinite
element.

\begin{result}{B.2.8}{Proved}{The first descent pullback is nonzero}
Let $D_P$ be the literal descent and put
$\chi_1=\Delta\circ D_P$.  There is an actual multiplier-one
gentle element $g$ with
\[
                       (\Delta(g),\chi_1(g))=(1,1).
                                                        \tag{B.N1}
\]
Together with the element $k$ of B.1.3 this gives
\[
 (\Delta(k),\chi_1(k))=(1,0),\qquad
 (\Delta(gk),\chi_1(gk))=(0,1).                \tag{B.N2}
\]
Both characters vanish on $\widehat{GT}$, and they are
linearly independent over $\mathbb F_2$.  In particular
$\Delta$ does not generate the continuous quadratic
annihilator of $\widehat{GT}$ in $\widehat{GT}_0$.
The subgroup $\ker\Delta$ is not preserved by $D_P$.
\end{result}

\begin{proof}
We first give the complete finite cover certificate,
including the proof that no central part was omitted.
Use $D$, $D^{\mathrm c}$ and $M$ from (B.M1)--(B.M3).
The exact calculations in
\texttt{check\_descent\_central\_full.g} give
\[
 M\cong C_2^3\times C_3,
 \qquad |D^{\mathrm c}|=9393093476352.           \tag{B.N3}
\]
Here is a reproducible finite justification of the full
kernel assertion, rather than just its first elementary
quotient.  For a central extension by $\mathbb F_p$,
evaluate a cocycle on the source torsion relation of
order $p$.  At a later step also evaluate the commutators
of every old central-kernel generator with every
polycyclic generator.  These evaluations are linear
functions of the cohomology class, independent of the
choices of lifts.  Their common kernel describes exactly
the central extensions that both receive the marked
$B_0$-map and remain central over the original $D$.
The relation of the other torsion generator can always
be corrected, since its order is prime to $p$.

The saved exact matrices have the following dimensions
and ranks.  Denote by $E_2$ and $E_3$ the maximal
elementary central two- and three-quotients of the
marked source over $D$.
\[
\begin{array}{c|c|c|c|c}
 \text{base}&p&\dim H^2(\text{base},\mathbb F_p)
       &\operatorname{rank}(\text{source/centrality conditions})
       &\dim\ker\\ \hline
 D&2&4&1&3\\
 D&3&2&1&1\\
 E_2&2&3&3&0\\
 E_3&3&1&1&0
\end{array}                                             \tag{B.N4}
\]
The first two rows construct $E_2$ and $E_3$, with
kernels $C_2^3$ and $C_3$.  They are actual quotients of
$B_0$, generated by their distinguished order-two and
order-three elements.  The last two rows prove that the
corresponding primary components of $M$ have exponents
two and three.  Indeed a nonzero $pM_p$ would have a
nonzero quotient $C_p$, producing a further marked
central extension of $E_p$ over $D$.  Its cohomology
class cannot be zero: a split extension would have
abelianization $(E_p)^{\mathrm{ab}}\times C_p$, whereas
every displayed marked source quotient has
abelianization $B_0^{\mathrm{ab}}\cong C_6$.
This contradicts the zero last kernel in (B.N4).
The transfer argument of B.2.6 shows that the primes
dividing $|M|$ divide $|D|=2^{29}3^6$, so there are no
other primary components.  This proves (B.N3).

The actual joint image of $B_0$ in $E_2\times_D E_3$
has kernel $C_2^3\times C_3$ over $D$: the two kernels
have coprime orders, and their joint image surjects
onto both.  The program verifies this joint image
has the order in (B.N3), is generated by the
distinguished source pair, and has central kernel
with invariant factors $[2,2,2,3]$.  Consequently it
is the full marked $D^{\mathrm c}$.

We now specify the finite automorphism used for (B.N1).
In the six degree-twenty blocks of the actual source
representation (B.M2), put
\[
 C=\begin{pmatrix}7&6\\0&7\end{pmatrix}
       \in\operatorname{SL}_2(\mathbb Z/12),
 \qquad (C,C,C,I,I,C),                           \tag{B.N5}
\]
and put the identity in the final modular block.
This is candidate fourteen in the executed source
calculation.  The code verifies that this tuple is
an actual element $f_D\in\overline F_0'$ and that
$\overline s\mapsto\overline s f_D$,
$\overline r\mapsto f_D^{-1}\overline r$ is an
automorphism of $D$.

In the full cover (B.N3), formula (B.M6) gives its
unique cusp-fixed lift $f^{\mathrm c}$.  The exact
finite checks give
\[
\begin{gathered}
 f^{\mathrm c}\in P',\qquad
 (S_0f^{\mathrm c})^2=((f^{\mathrm c})^{-1}R_0)^3=1,\\
 S_0f^{\mathrm c}=\beta^{\mathrm c}S_0
                     (\beta^{\mathrm c})^{-1},\qquad
 (f^{\mathrm c})^{-1}R_0=A^{\mathrm c}R_0
                     (A^{\mathrm c})^{-1},\\
 \beta^{\mathrm c},A^{\mathrm c}\in P,
 \qquad q_{\mathrm c}(\beta^{\mathrm c})
       =q_{\mathrm c}(A^{\mathrm c})=[5I].        \tag{B.N6}
\end{gathered}
\]
The resulting endomorphism of $D^{\mathrm c}$ is
bijective, both by B.2.6 and by the independently
executed finite check.  The certificate file retains
the complete polycyclic presentations, the cocycles
and centrality matrices of (B.N4), and the literal
polycyclic coordinates of $f^{\mathrm c}$,
$\beta^{\mathrm c}$ and $A^{\mathrm c}$ in (B.N6).
Thus the two required positive frames are actual
elements, not only asserted conjugacy-class labels.

Apply the image equality (B.M12) to this finite
automorphism and these two frames.  It supplies a
full positive automorphism $\Phi$ of $\widehat B_0$
fixing $\tau$, with the stated $D$-action and frames.
Write $f=s\Phi(s)$.  The first block and the
$r^2$ block of (B.N5) give, in the actual literal
pants quotient,
\[
 q\bigl(f(s,\tau)\bigr)=[C],\qquad
 q\bigl(f(r^{-1},s)\bigr)=1,
 \qquad q(\beta)=q(A)=[5I].                    \tag{B.N7}
\]
The block identification follows from the exact
source conjugation table
$r(x,y)r^{-1}=(y,z)$ and
$r^2(x,y)r^{-2}=(z,x)$, with
$\varphi_P(x,y,z)=(s,\tau,r^{-1})$.
The signed matrix for the third peripheral can
differ by $-I$, which is killed by the actual map
to $Q$; thus (B.N7) holds before imposing any
commutator normalization on $f$.

We first normalize $\Phi$ and compute its own
character.  Write $\beta^{\mathrm{ab}}=(u,v)$,
$a=v-u$.  Formula (B.M13) gives
$\Phi_0=\operatorname{Inn}(x^a)\Phi$ and a full
gentle element $g=(1,f_0)$ with frames
$\beta_0=x^a\beta$, $A_0=x^a A$.
Its cubic frame has zero pure abelianization.
To see this directly, the normalized commutator
condition gives
$(1-\operatorname{Ad}_r)A_0^{\mathrm{ab}}=0$;
the matrix of $1-\operatorname{Ad}_r$ has determinant
three, and multiplication by three is injective
on $\widehat{\mathbb Z}^{2}$.  Hence
$A^{\mathrm{ab}}=(-a,0)$.

In the canonical identification
$q(F)=V_4\times\operatorname{SL}_2(\mathbb F_3)$,
the residue $[5I]$ has coordinates $(I,-I)$.
The $V_4$ factor of $q(A)$ implies $a\equiv0\pmod2$.
The quotient
$\operatorname{SL}_2(\mathbb F_3)\to C_3$
sends both pure generators $x,y$ to the same
generator; applied to $A$ it gives $a\equiv0\pmod3$.
Applied to $\beta$, these same two maps give
$u,v\equiv0\pmod2$ and $u+v\equiv0\pmod3$.
Since $v-u=a\equiv0\pmod3$, both $u$ and $v$
are divisible by six.  Therefore
\[
 q(\beta_0)=q(A_0)=[5I],\qquad
 \rho_g=-v,\qquad q(V_g)=U^{8v}=I.             \tag{B.N8}
\]
The two signed CM bits are consequently $(1,0)$,
as in the exact $[5I]$ computation of B.1.3, and
$\Delta(g)=1$.

It remains to compute the descendant, while retaining
the inner correction just made.  Put
$C_f=f(s,\tau)$ and $H_f=f(r^{-1},s)$ in the
full target $\widehat B_0$.  Their residues in
(B.N7) are trivial modulo two, so both are pure.
Write $C_f^{\mathrm{ab}}=(p,q_0)$.  We justify the
raw formulas without assuming $f\in\widehat F_0'$.
The actual source torsion relations imply the two
cycle identities for this $f$.  Using
$r(x,y)r^{-1}=(y,z)$ and $\Phi(r)=f^{-1}r$, they give
\[
\begin{split}
 \Phi(z)&=\Phi(ryr^{-1})
       =f^{-1}f(y,z)^{-1}z f(y,z)f\\
       &=f(z,x)z f(z,x)^{-1}.
\end{split}
\]
Thus $\Phi|_{\widehat F_0}$ preserves the closed
normal subgroup generated by $x^2,z^3$, in both
directions, and induces an actual automorphism of
its quotient $\varphi_P$.  It fixes $s$, conjugates
$\tau$ by $C_f^{-1}$, and conjugates $r$ by $H_f$.
After conjugation by $C_f$ its pair coordinate is
$\sigma(C_f)C_f^{-1}$, with abelianization
$(q_0-p,p-q_0)$.  The further conjugation by
$x^{q_0-p}$ makes this coordinate a commutator.
It gives the exact target frames
\[
 \beta'=x^{q_0-p}C_f,\qquad
 A'=x^{q_0-p}C_fH_f,\qquad
 \rho'=-q_0.                                   \tag{B.N9}
\]
The zero abelianization of $A'$ gives
$H_f^{\mathrm{ab}}=(-q_0,-q_0)$.

The $V_4$-normalized representative of $[C]$
has modulo-four value $U^2$ and modulo-three
value $-I$: multiply the displayed matrix in
(B.N5) by its common sign $-1$.  Thus the
$V_4$ component of (B.N7) gives $p$ odd and
$q_0$ even.  The $C_3$ quotient just used gives
$p+q_0\equiv0\pmod3$ from $C_f$, and
$-2q_0\equiv0\pmod3$ from $H_f$.  Consequently
\[
 q_0\equiv0\pmod6,\qquad p\equiv3\pmod6,
 \qquad q_0-p\equiv3\pmod6.                   \tag{B.N10}
\]
Since $U^6C=7I$ modulo twelve, equations
(B.N9)--(B.N10) give
\[
 q(\beta')=q(A')=[7I]=[5I],\qquad
 q(V')=U^{8q_0}=I.                              \tag{B.N11}
\]
This descendant again has the two signed CM bits
$(1,0)$, hence character value one.

Finally the source correction
$\Phi_0=\operatorname{Inn}(x^a)\Phi$ induces
$\operatorname{Inn}(\varphi_P(x)^a)
=\operatorname{Inn}(s^a)$ on the descendant.
It preserves its outer class; here $a$ is divisible
by six, so $s^a=1$ and the two raw quotient
automorphisms even coincide literally.  The
faithful normalized identification explained after
(B.M13) makes the normalized descendant the same
gentle element.  Therefore the value in
(B.N11) is exactly $\Delta(D_P(g))$, proving
(B.N1) for the actual normalized full gentle
element.  B.1.8 gives $\chi_1|_{\widehat{GT}}=0$,
while A.7 gives $\Delta|_{\widehat{GT}}=0$.
The actual kernel element $k$ has
$\Delta(k)=1$ and $D_P(k)=1$, so (B.N2) follows.
These two value vectors prove independence and
the asserted failure of preservation of
$\ker\Delta$ under $D_P$.
\end{proof}
